% Generated from paper_v1.md; original Markdown is unchanged.
\documentclass{article}
% Anonymous workshop submission. Do not use final or preprint for review.
\PassOptionsToPackage{numbers,sort&compress}{natbib}
\usepackage[dblblindworkshop]{neurips_2026}
\workshoptitle{Physical World AI: Geometry, Characteristics, and Multimodal Sensing}
\usepackage[T1]{fontenc}
\usepackage{amsmath,amssymb}
\usepackage{graphicx,booktabs,longtable,array,calc}
\usepackage{microtype,xcolor,newunicodechar}
\usepackage{url}
\usepackage[hidelinks,unicode]{hyperref}
\newunicodechar{−}{\ensuremath{-}}
\newunicodechar{±}{\ensuremath{\pm}}
\newunicodechar{×}{\ensuremath{\times}}
\newunicodechar{Δ}{\ensuremath{\Delta}}
\newunicodechar{θ}{\ensuremath{\theta}}
\newunicodechar{φ}{\ensuremath{\phi}}
\newunicodechar{ε}{\ensuremath{\epsilon}}
\newunicodechar{→}{\ensuremath{\rightarrow}}
\newunicodechar{≤}{\ensuremath{\leq}}
\newunicodechar{≥}{\ensuremath{\geq}}
\providecommand{\tightlist}{\setlength{\itemsep}{0pt}\setlength{\parskip}{0pt}}
\providecommand{\pandocbounded}[1]{#1}
\setlength{\emergencystretch}{2em}
\graphicspath{{figures/}{../figures/}}
\hypersetup{pdftitle={Auditing Gait Laterality in Skeleton Predictive Representations},pdfauthor={Anonymous Author(s)},pdfcreator={LaTeX},pdfsubject={Physical World AI workshop submission}}
\title{Auditing Gait Laterality in Skeleton Predictive Representations}
\author{Anonymous Author(s)}
\begin{document}
\maketitle
\begin{abstract}
Human movement provides a useful test of whether predictive representations preserve the geometry of an articulated body. Left--right reflection supplies a known transformation, while differences between the two sides provide an observable quantity whose sign should change under that transformation. We study this relationship in a skeleton Joint-Embedding Predictive Architecture using pose sequences from gait videos. A source-separated evaluation compares trained and initial encoders, tests reflection behavior, and changes the anatomical and temporal placement of hidden training targets. The available experiments show that predicting a clip's hidden features more accurately than another clip's features can coexist with weaker recovery of a signed movement contrast. Motion-sensitive summaries improve access to information already present at initialization, while motion-weighted and connected-region masks provide no demonstrated advantage over their matched random references. We distinguish this evidence from exact symmetry imposed at the output and from synthetic demonstrations of explicit reflection training. The study contributes an evaluation of articulated-body representations for Physical World AI, identifying a gap between feature prediction and the preservation of a task-relevant physical measurement.
\end{abstract}
\section{Motivation and research question}\label{motivation-and-research-question}

Walking involves coordinated movement of an articulated body with corresponding left and right limbs. This geometry provides both useful regularities and meaningful departures from them. An average description of movement can remain unchanged when anatomical sides are exchanged, whereas a quantity describing which side moves more should reverse sign. A representation intended for physical prediction may need to support both responses.

Gait abnormalities motivate studying that distinction. Post-stroke studies measure spatial and temporal asymmetries, and Parkinson's research examines unequal arm swing Patterson et al.~\citep{ref01}, Lewek et al.~\citep{ref02}. These observations motivate our measurement; they do not validate it as a clinical index. Here, laterality means a signed comparison of estimated left- and right-side movement.

We ask whether masked feature prediction makes that comparison easier to recover from a skeleton encoder on held-out source videos. The initial null is that training provides no improvement over the same encoder at initialization under a fixed downstream evaluation. Separate questions concern whether the output changes sign correctly and whether landmark features follow a specified reflection rule. Answering one does not establish the others.

Our contribution is a controlled evaluation of a small predictive representation model. The completed experiments concern whole-clip masked prediction and frozen-feature readouts. Their relevance to world models is the assessment of a geometric property that could matter to a future dynamics model. The study has no demonstrated real-data forecasting, planning, or multimodal-fusion result.

\section{Data and the geometric target}\label{data-and-the-geometric-target}

The local subset of the Gait Abnormality in Video Dataset~\citep{ref03} contains 666 annotations, of which 642 have pose archives. Quality control retains 625 sequences from 93 source videos. The annotations cover normal, Parkinson's, stroke, myopathic, and cerebral-palsy gait. They stratify source assignments and describe the dataset; they do not supervise encoder training.

Each archive contains a variable-length sequence of 33 landmarks with estimated coordinates and visibility. Coordinates are derived from video, including an inferred depth component, and are not calibrated metric three-dimensional motion capture.

For each of five bilateral pairs---shoulders, knees, ankles, heels, and foot tips---we calculate left and right speeds using observed timestamps and only transitions for which both landmarks are observed at both ends. Let their median speeds be \(m_{L,k}\) and \(m_{R,k}\). The target is

\[
y(x)=\frac{1}{5}\sum_{k=1}^{5}
\frac{m_{L,k}-m_{R,k}}{m_{L,k}+m_{R,k}+\epsilon}.
\]

All five pairs must provide at least eight shared valid transitions. Hips establish the pelvis reference but are excluded from this target: pelvis centering makes their speed magnitudes identical. The target is calculated before input interpolation and resizing.

Reflection \(M\) negates the centered horizontal coordinate and exchanges all anatomically paired landmarks, together with validity. Consequently, \(M^2x=x\) and \(y(Mx)=-y(x)\). This identity is a property of the coordinate operation and formula, rather than evidence that every reflected recording would occur naturally.

The encoder input is prepared separately. Short gaps can be interpolated, coordinates are normalized, and the sequence is resampled to \(64\times33\times3\). A matching validity array accompanies it. Four time steps form one landmark token, producing a grid of 16 time blocks by 33 landmarks. Anatomical identity and array structure are preserved; the original trajectories, timestamps, and motion magnitudes are not preserved exactly by resampling.

\section{Training and source separation}\label{training-and-source-separation}

All clips and derived views from a source video inherit its fold. Five outer folds contain respectively 19, 19, 19, 18, and 18 test sources; their test sequence counts are 189, 182, 72, 77, and 105. Each encoder trains only on its matching outer-training sources. New encoders are initialized for each fold and seed, and source-balanced sampling limits domination by videos with many clips.

The base model uses a 96-dimensional embedding, four encoder layers, two predictor layers, and four attention heads. An online encoder processes visible tokens; a predictor estimates hidden teacher features. A slowly updated target encoder supplies those features without receiving gradients. A variance and covariance regularizer encourages features to retain variation. The base run uses five seeds and two reflection variants, with a recorded budget of 1,200 optimizer updates per encoder.

Laterality enters the completed training recipe through the named anatomical schema, the eligible landmark set, and a comparison that reflects training inputs with probability one half. These choices supply anatomical structure without supplying the signed target. Notebook 09 additionally implements a penalty aligning features of original and reflected inputs after anatomical exchange. Its retained execution is synthetic, so the penalty has no completed real-GAVD efficacy result.

After pretraining, encoder weights are frozen. Inner source-separated validation selects a ridge-regression penalty; ridge regression is a linear readout whose weight penalty limits fitting unstable feature combinations. The original audit uses four inner folds, while later comparisons use three. Imputation and scaling are fitted within each inner fitting partition and refitted on the outer-training set after selection. These inner folds select the readout only: the encoder has already seen their unlabeled inputs during outer-training pretraining.

\section{Evaluation through successive questions}\label{evaluation-through-successive-questions}

\subsection{Does the original representation satisfy the reflection audit?}\label{does-the-original-representation-satisfy-the-reflection-audit}

The retained original aggregate summary reports source-balanced \(R^2=0.060\), with a 95\% source-bootstrap interval of \([-0.025,0.126]\). The learned-minus-initial difference is \(-0.018\), with interval \([-0.039,0.002]\). Reflection augmentation changes \(R^2\) by \(0.004\), with interval \([-0.006,0.013]\). These estimates do not demonstrate a training benefit or establish equivalence between variants.

The direct token test aligns landmark positions under reflection while leaving feature channels unchanged. Its normalized squared discrepancy is

\[
q=\frac{\|Z(Mx)-SZ(x)\|_C^2}
{\|Z(Mx)\|_C^2+\|SZ(x)\|_C^2},
\]

where \(S\) exchanges joint identities and \(C\) selects common valid tokens. The retained summary reports \(q=0.114\) for learned features and \(0.083\) at initialization, with a paired increase of \(0.031\), interval \([0.016,0.048]\). This is unfavorable evidence for the specified identity-channel transformation. An encoder could express reflection through a different transformation of feature channels; the test does not examine that possibility.

An output can be made antisymmetric through an odd feature, \(z^-(x)=[z(x)-z(Mx)]/\sqrt{2}\), followed by a linear readout without an intercept or feature centering. Its symmetry follows algebraically. Predictive accuracy still requires evidence: the recorded learned-minus-initial difference under this construction is negative.

The original values are available in a retained figure-summary JSON. Its full report, checkpoint, and per-example prediction chain was not found in this checkout during the present review. We retain these results as previously reported context and distinguish them from the newly audited comparisons below.

\subsection{Does anatomical target selection improve learning?}\label{does-anatomical-target-selection-improve-learning}

Notebook 12 compares scattered hidden targets selected from twelve gait landmarks with the same number selected across all 33 landmarks. Both encoders still receive the full landmark set and retain anatomical choices in regularization and readout. Starting weights, sampled clips, geometric views, and 1,200 updates are paired.

The teacher-feature readouts obtain \(R^2=-0.023\) with gait targets and \(-0.010\) with all-landmark targets; the initial encoder obtains \(0.048\). The broader-minus-gait difference has interval \([-0.018,0.050]\), leaving both improvement and harm compatible with these recordings. Larger ridge penalties frequently reach the edge of the search range, limiting conclusions about the best achievable readout.

\subsection{Does targeting motion or connected anatomy help?}\label{does-targeting-motion-or-connected-anatomy-help}

Notebooks 15--18 extend the comparison to motion-weighted and connected-region targets. The completed grid contains 125 trained encoders, each with 1,200 updates. Motion comparisons and region comparisons have their own random references, matched to the actual hidden count. Equal nominal mask percentages would not ensure equal tasks because missing observations and eligible landmark pools differ.

\begin{table}[htbp]
\centering
\small
\begin{tabular}{@{}lr@{}}
\toprule\noalign{}
Frozen features or control & Mean source-balanced \(R^2\) \\
\midrule\noalign{}

Training-source target mean & −0.0106 \\
Initial encoder, mean summary & 0.0708 \\
Initial encoder, motion-sensitive summary & 0.2225 \\
Trained teachers, motion-sensitive summaries & 0.1008--0.1142 \\
\bottomrule
\end{tabular}
\end{table}

The motion-sensitive summary adds temporal variation, feature changes, and observation support. All trained arms remain below their matched initial encoder in each of the five seeds. Its advantage over mean pooling cannot yet be attributed solely to motion because the summary also includes support information.

\begin{table}[htbp]
\centering
\small
\begin{tabular}{@{}
  >{\raggedright\arraybackslash}p{(\linewidth - 4\tabcolsep) * \real{0.4600}}
  >{\raggedleft\arraybackslash}p{(\linewidth - 4\tabcolsep) * \real{0.2200}}
  >{\raggedright\arraybackslash}p{(\linewidth - 4\tabcolsep) * \real{0.3200}}@{}}
\toprule\noalign{}
\begin{minipage}[b]{\linewidth}\raggedright
Mask minus its matched random reference
\end{minipage} & \begin{minipage}[b]{\linewidth}\raggedleft
\(\Delta R^2\)
\end{minipage} & \begin{minipage}[b]{\linewidth}\raggedright
95\% source interval
\end{minipage} \\
\midrule\noalign{}

MAMP-style motion weighting & −0.0008 & {[}−0.0204, 0.0153{]} \\
Robust motion mixture & 0.0005 & {[}−0.0155, 0.0152{]} \\
Connected region & −0.0087 & {[}−0.0333, 0.0177{]} \\
\bottomrule
\end{tabular}
\end{table}

The mask audits confirm that motion weighting changes target selection and connected regions remove local temporal clues. Their intended geometric effects therefore occurred, but no alternative mask demonstrates improved target prediction.

\subsection{What does successful feature prediction establish?}\label{what-does-successful-feature-prediction-establish}

The trained predictors consistently predict their own clip's hidden teacher features more accurately than targets from another source video under the recorded diagnostic. Initial predictors show no consistent preference. This establishes useful clip correspondence for that diagnostic. Posture, viewpoint, observation support, or motion could contribute to it; the result does not isolate a physical mechanism.

Thus, predictor training can succeed on its own features while the frozen representation yields weaker laterality readouts. Because each arm supplies its own changing teacher, raw feature-prediction losses are not a common between-model measurement scale. The coordinate-derived outcome supplies a shared evaluation target.

\section{Statistical interpretation}\label{statistical-interpretation}

Each source receives equal total weight. Within a seed, predictions from all five held-out folds are pooled before computing \(R^2\); the reported result averages the five seed-specific scores. It is neither an average of fold \(R^2\) values nor a score from an ensemble of seed-averaged predictions. Negative \(R^2\) means squared error exceeds variation around the weighted evaluation mean; a separately fitted training-source mean is the usable prediction baseline.

Paired uncertainty resamples entire videos while retaining their clips and paired seed predictions. The 2,000-resample intervals condition on the fitted models and exclude full-retraining and new-split uncertainty. Repeated seeds and diagnostic mask draws are not independent participants. Subsequent experiments were informed by the same development cohort, so their source separation does not make the overall research trajectory an untouched confirmatory study.

\section{Implications and limitations}\label{implications-and-limitations}

The strongest supported conclusion concerns access to one signed movement measurement under the evaluated readouts. A negative learned-over-initial contrast does not demonstrate that all motion information has disappeared. Mean pooling, feature scaling, regularization, and the difference between the original target and prepared input remain possible contributors.

The study fits the workshop's articulated geometry and evaluation themes. Coordinates derived from RGB do not constitute multimodal sensing, and an inferred depth coordinate does not establish calibrated 3D reconstruction. Clinical interpretation would require independent measurements, affected-side labels, and participant-separated validation. The source-video split cannot exclude the same unidentified person appearing in different videos.

The next experiments should first isolate support features, broaden training-only ridge selection, and measure the target after each preparation stage. A later study can test whether explicit geometric training improves common observable future endpoints, using past-only preprocessing and persistence and velocity baselines. Existing symmetry-training and future-decoding demonstrations are synthetic; they are not evidence of forecasting gait recordings.

Institutional ethics, data-use, and derived-pose release reviews remain unresolved in the recorded governance status. This manuscript is an internal revision, and it does not claim those determinations have been obtained.

\section{Related work and reproducibility}\label{related-work-and-reproducibility}

S-JEPA~\citep{ref04} already studies masked skeleton-feature prediction, and MAMP~\citep{ref05} motivates motion-weighted masking. SLiM~\citep{ref06} provides recent structured skeleton-masking precedent. Our contribution is the controlled laterality evaluation, with no claim to invent skeleton JEPA or anatomical masking.

The {tutorial}, {source-split reference}, {core audit}, and {extension audit} document the evidence and its limits. Original notebooks and experiment artifacts are preserved. This revision adds no training run or clinical finding.
\label{physworld:main-last}
\clearpage
\label{physworld:references-start}
{\renewcommand{\bibfont}{\small}
\begin{thebibliography}{99}
\bibitem{ref01} K. K. Patterson, I. Parafianowicz, C. J. Danells, V. Closson, M. C. Verrier, W. R. Staines, S. E. Black, and W. E. McIlroy. \href{https://pubmed.ncbi.nlm.nih.gov/18226655/}{Gait asymmetry in community-ambulating stroke survivors}. \emph{Archives of Physical Medicine and Rehabilitation}, 89(2):304--310, 2008. doi:10.1016/j.apmr.2007.08.142.

\bibitem{ref02} M. D. Lewek, R. Poole, J. Johnson, O. Halawa, and X. Huang. \href{https://pubmed.ncbi.nlm.nih.gov/19945285/}{Arm swing magnitude and asymmetry during gait in the early stages of Parkinson's disease}. \emph{Gait \& Posture}, 31(2):256--260, 2010. doi:10.1016/j.gaitpost.2009.10.013.

\bibitem{ref03} R. Ranjan, D. Ahmedt-Aristizabal, M. A. Armin, and J. Kim. \href{https://arxiv.org/abs/2407.04190}{Computer Vision for Clinical Gait Analysis: A Gait Abnormality Video Dataset}. \emph{IEEE Access}, 13:45321--45339, 2025. doi:10.1109/ACCESS.2025.3545787.

\bibitem{ref04} M. Abdelfattah and A. Alahi. \href{https://www.ecva.net/papers/eccv_2024/papers_ECCV/papers/04755.pdf}{S-JEPA: A Joint Embedding Predictive Architecture for Skeletal Action Recognition}. In \emph{Computer Vision -- ECCV 2024}, volume 15090 of \emph{Lecture Notes in Computer Science}, pages 367--384. Springer, 2025. doi:10.1007/978-3-031-73411-3\_21.

\bibitem{ref05} Y. Mao, J. Deng, W. Zhou, Y. Fang, W. Ouyang, and H. Li. \href{https://arxiv.org/abs/2308.07092}{Masked Motion Predictors are Strong 3D Action Representation Learners}. In \emph{Proceedings of the IEEE/CVF International Conference on Computer Vision (ICCV)}, pages 10181--10191, 2023.

\bibitem{ref06} J. Do, Y. Chen, G. Youk, and M. Kim. \href{https://arxiv.org/abs/2603.10648}{Less is More: Compact-Token Masked Feature Prediction for Skeleton Representation Learning}. arXiv preprint arXiv:2603.10648v3, 2026.
\end{thebibliography}
}

\end{document}
